% !TEX root = ../main.tex
\chapter{Elektrisk integration og wiring}
\label{ch:el}

Dette kapitel beskriver den elektriske opbygning af det arm-delsystem,
projektet omfatter: hvilke komponenter der indgår, hvordan de er
forbundet, og hvilke integrationsproblemer der opstod undervejs.

\section{Komponentoversigt}

Det elektriske delsystem omfatter komponenter på tre niveauer: aktuatorer
i armen, en kommunikations-mikrocontroller mellem arm og host, og en host
der kører ROS\,2-laget. Komponenterne er listet i tabellen nedenfor
(se Bilag, Tabel~3). De væsentligste er:

\begin{itemize}
    \item 7~stk.\ Waveshare ST3215 seriebus-servoer, én pr.~led i venstre arm.
          % TODO ALEX: bekræft eller ret. Jeg har skrevet:
          % "indbygget magnetisk encoder og intern PID-regulator,
          % kommunikerer på en halv-duplex TTL-bus".
          % Spørgsmål: er det rigtigt? Er der andre spec-detaljer
          % du gerne vil have med (spænding, max moment, opløsning)?
          Servoerne har indbygget magnetisk encoder og intern PID-regulator
          og kommunikerer på en halv-duplex TTL-bus.
    \item 1~stk.\ ESP32-baseret kontrolboks (RØD), der konverterer mellem
          USB på host-siden og servobussens TTL-niveau. To tilsvarende
          bokse (GUL og HVID) findes til højre arm + hoved henholdsvis
          hænder, men ligger uden for projektets afgrænsning.
    \item 1~stk.\ NVIDIA Jetson Orin Nano, der kører ROS\,2 Humble og
          afvikler manager-noden samt commissioning-værktøjerne.
    \item Strømforsyning til servobussen samt logik-forsyning til
          ESP32-boksen via USB.
\end{itemize}

\section{Signalkæde og wiring}

Signalvejen i den oprindelige konstruktion er kort og lineær:

\begin{enumerate}
    \item Jetson Orin Nano sender positionskommandoer på et ROS-topic.
    \item Manager-noden serialiserer kommandoerne og skriver dem på en
          virtuel COM-port mod ESP32-boksen.
    \item ESP32-boksen konverterer beskeden til servobus-protokollen og
          videresender den på den halv-duplexerede TTL-bus.
    \item Hver ST3215-servo på bussen filtrerer på sit servo-ID og
          udfører den modtagne positionskommando i sin egen interne
          PID-sløjfe.
\end{enumerate}

Effektvejen er adskilt fra signalvejen: servobussen forsynes direkte fra
sin egen forsyning, mens ESP32 og logik-niveau strømmes via USB fra
Jetson'en.

\section{Stabil portbinding for ESP32-bokse}
\label{sec:usb-binding}

Den væsentligste elektriske/integrationsmæssige rettelse i projektet
vedrører tildelingen af USB-porte til de tre ESP32-bokse. Tidligere
blev de adresseret som \enquote{/dev/ttyUSB0}, \enquote{/dev/ttyUSB1}
og \enquote{/dev/ttyUSB2}, og rækkefølgen afhang af, hvilken boks
operativsystemet enumererede først ved boot. I praksis betød det, at
manager-noden kunne starte med at sende den venstre arms kommandoer
til den højre arms ESP32, hvis kablerne var sat i en anden rækkefølge. 
Systemet kører, det kører bare mod den forkerte hardware.

Løsningen er at adressere boksene via deres faste sti i
\enquote{/dev/serial/by-path/} frem for en bestemt tilslutningsrækkefølge
\enquote{/dev/ttyUSBx}. Hver fysisk USB-port på Jetson-hubben har en
unik sti, der ikke ændrer sig mellem boots eller hvis maskinen
genstartes. Manager-nodens konfiguration blev opdateret til at slå op
i denne sti, og der blev oprettet dokumentation for hvor hver ESP32 sluttes til.

\begin{itemize}
    \item Port 2.2 -- RØDT kabel -- venstre arm
    \item Port 2.3 -- GULT kabel -- højre arm og hoved
    \item Port 2.1 -- HVIDT kabel -- hænder
\end{itemize}

Mappingen afhænger nu af, hvilken fysisk USB-port boksen sidder i,
ikke af om operatøren husker rækkefølgen ved tilslutning.

\section{Praktiske valg under integrationen}

% TODO ALEX: Slet hele denne sektion hvis der ikke er noget faktisk
% at sige om kabelføring/stik/mekanisk mærkning, der adskiller
% sig fra det oprindelige design.
% Hvis der ER noget, fx:
%   - Har du ændret på kabelføringen?
%   - Er der specifikke trækaflastninger du har lavet?
%   - Hvordan har du mærket stik/porte fysisk? (du nævnte selv
%     nummerering på din replica frem for farvekode)
% Skriv det her, så omskriver jeg.

% TODO: konkret strømbudget (max strøm pr. servo x antal samtidigt
% aktive + dimensionering af forsyning) når data foreligger fra lab.

